%=====================================================================
% AGPS — IEEE Conference Paper
% Template: IEEEtran (official IEEE conference format)
%
% Compile on Overleaf:  New Project -> Upload Project -> this file
% Or locally:           pdflatex -> bibtex -> pdflatex -> pdflatex
%=====================================================================

\documentclass[conference]{IEEEtran}
\IEEEoverridecommandlockouts

\usepackage{cite}
\usepackage{amsmath,amssymb,amsfonts}
\usepackage{algorithmic}
\usepackage{graphicx}
\usepackage{textcomp}
\usepackage{xcolor}
\usepackage{url}
\usepackage{booktabs}

\def\BibTeX{{\rm B\kern-.05em{\sc i\kern-.025em b}\kern-.08em
    T\kern-.1667em\lower.7ex\hbox{E}\kern-.125emX}}

\begin{document}

%---------------------------------------------------------------------
\title{A Deterministic Rule-Based Multi-Criteria Framework for
Transparent and Explainable Government Tender Evaluation}
% Keep the title specific. Avoid "novel", "revolutionary", "AI-powered".

% ---------------------------------------------------------------------
% LAYOUT A -- all authors share one institution (use this one).
% IEEEtran places at most 3 author blocks per row, so 5 separate blocks
% overflow the column. Grouping the names under a single shared
% affiliation is the standard fix and is accepted by IEEE.
% ---------------------------------------------------------------------
\author{
\IEEEauthorblockN{
Manav Singh,
Second Author,
Third Author,
Fourth Author,
Fifth Author}
\IEEEauthorblockA{\textit{Dept. of Computer Science and Engineering} \\
\textit{Your Institute Name}, City, India \\
manav@example.com, second@example.com, third@example.com, \\
fourth@example.com, fifth@example.com}
}

% ---------------------------------------------------------------------
% LAYOUT B -- authors from DIFFERENT institutions, or if you want each
% author listed with their own affiliation block. Delete Layout A above
% and uncomment this. Three blocks per row; \\ starts a new row.
% ---------------------------------------------------------------------
%\author{
%\IEEEauthorblockN{Manav Singh}
%\IEEEauthorblockA{\textit{Dept. of CSE} \\ \textit{Institute Name} \\
%City, India \\ manav@example.com}
%\and
%\IEEEauthorblockN{Second Author}
%\IEEEauthorblockA{\textit{Dept. of CSE} \\ \textit{Institute Name} \\
%City, India \\ second@example.com}
%\and
%\IEEEauthorblockN{Third Author}
%\IEEEauthorblockA{\textit{Dept. of CSE} \\ \textit{Institute Name} \\
%City, India \\ third@example.com}
%\\
%\IEEEauthorblockN{Fourth Author}
%\IEEEauthorblockA{\textit{Dept. of CSE} \\ \textit{Institute Name} \\
%City, India \\ fourth@example.com}
%\and
%\IEEEauthorblockN{Fifth Author}
%\IEEEauthorblockA{\textit{Dept. of CSE} \\ \textit{Institute Name} \\
%City, India \\ fifth@example.com}
%}

\maketitle

%---------------------------------------------------------------------
\begin{abstract}
% 150-250 words. No citations. Structure:
% (1) problem  (2) approach  (3) what was built  (4) key result
% (5) principal limitation
Public procurement evaluation is frequently opaque \dots
\end{abstract}

\begin{IEEEkeywords}
audit trail, decision support systems, explainable systems,
multi-criteria decision making, public procurement, simple additive
weighting
\end{IEEEkeywords}

%---------------------------------------------------------------------
\section{Introduction}
\label{sec:intro}

% Problem context: opacity and inconsistency in manual tender evaluation.
% Motivation: reproducibility and explainability as governance requirements.
% Objectives.

The contributions of this work are:

\begin{itemize}
\item A two-layer evaluation model combining non-compensatory rule-based
      screening with compensatory weighted aggregation.
\item A configuration-snapshot mechanism providing verifiable
      reproducibility of historical evaluations.
\item An explicit data-provenance model that surfaces the reliability of
      each input rather than absorbing it into the arithmetic.
\item A taxonomy of criterion verifiability distinguishing
      registry-verifiable attributes from forward-looking commitments.
\end{itemize}

The remainder of this paper is organised as follows.
Section~\ref{sec:related} reviews related work \dots

%---------------------------------------------------------------------
\section{Related Work}
\label{sec:related}

% MCDM for supplier selection: SAW, AHP, TOPSIS.
% Fuzzy / grey methods for uncertain input data.
% E-procurement transparency.
%
% Positioning statement (important):
% Existing approaches address input uncertainty through fuzzy or grey
% modelling, at the cost of exact reproducibility. This work retains
% crisp determinism and instead makes data provenance explicit.

%---------------------------------------------------------------------
\section{System Architecture}
\label{sec:arch}

AGPS was created as a rule-based procurement evaluation system where the
decision process is split into stages. The architecture separates
eligibility screening from the later ranking of eligible bids. This split
stops a bidder from making up for a missing requirement by doing well in
other evaluation criteria. I find this design clear and fair.

\subsection{Two-Layer Decision Model}
\label{subsec:two-layer}

The evaluation process has two layers. The first layer does eligibility
screening with procurement rules. Each rule is checked separately for
every submitted bid. If a bid breaks one or more rules it is marked as
ineligible. The screening does not stop after the failure; instead all
failed rules are gathered and returned so that the bidder and the
procuring entity can see all eligibility problems.

The eligibility layer is therefore non-compensatory. A bidder that fails
a requirement cannot regain eligibility by scoring higher on price,
quality, delivery, experience or any other criterion.

The second layer ranks bids that have passed the eligibility screening.
This layer uses Simple Additive Weighting (SAW). Each criterion is put on
a scale and multiplied by its set weight. The weighted values are added
together to give the score. I appreciate how the ranking layer lets
strong points balance ones.

The ranking layer is compensatory because a strong result on one
criterion can balance a result on another based on the declared weights.
This design matches procurement rules where mandatory conditions are kept
separate from the comparison of financially acceptable bids
\cite{ref1}, \cite{ref2}.
% [CITE ref1: Government of India, General Financial Rules 2017]
% [CITE ref2: CVC two-envelope evaluation practice]

\begin{figure}[htbp]
\centerline{\includegraphics[width=\columnwidth]{fig1_workflow.pdf}}
\caption{AGPS two-layer evaluation workflow showing bid submission
$\rightarrow$ eligibility screening $\rightarrow$ eligible bids
$\rightarrow$ criterion normalization $\rightarrow$ weighted aggregation
$\rightarrow$ deterministic ranking $\rightarrow$ award decision.}
\label{fig:workflow}
\end{figure}

\subsection{Generic Criterion Architecture}
\label{subsec:generic}

The scoring engine was built to work without being tied to procurement
criteria. Criteria are set up as configuration of being hard-coded into
the evaluation logic. Each criterion lists its preference direction,
weight and source of its value.

The system supports four types of value sources: \texttt{BID\_FIELD},
\texttt{VENDOR\_FIELD}, \texttt{TECHNICAL\_VALUE} and
\texttt{DERIVED\_QUALITY}. This lets criteria get values from a bid, from
the vendor profile, from technical evaluation data or from a
deterministic quality-derivation process.

The default setup has four criteria: price at 40\%, quality at 30\%,
delivery at 20\% and experience at 10\%. Price and delivery are
lower-is-better while quality and experience are higher-is-better.

Because criteria are driven by configuration the evaluation engine does
not have logic for each procurement attribute. A demo tender with six
criteria, such as Warranty and Maintenance Support, was evaluated without
changing the evaluation engine. This design lets the same calculation
process be used for tenders that have different evaluation structures.

\subsection{Configuration Snapshot and Graduated Locking}
\label{subsec:locking}

To stop evaluation rules from changing after bidders have made their
submissions AGPS uses a graduated locking mechanism. A tender starts in
the \texttt{UNLOCKED} state, where its configuration can be edited
freely.

After publication but before any bid arrives the tender moves to the
\texttt{SOFT\_LOCKED} state. In this state changes can still happen. They
are versioned and logged in the audit trail.

The first bid that comes in causes a move to the \texttt{HARD\_LOCKED}
state. From then on the tender configuration cannot change. Big changes
after bids are in place need the tender to be cancelled and re-issued,
not just the evaluation rules altered.

This system makes it clear that no bidder is judged by rules that were
changed after that bidder made a bid.

\subsection{Determinism Contract}
\label{subsec:determinism}

AGPS uses ways to make the evaluation repeatable. Before bidding the full
tender configuration is taken as a snapshot. A \texttt{configHash} is
then computed with SHA-256 over the key-sorted JSON version of that
snapshot.

The ranking process adds values in a fixed order of criteria. The
calculation is done with float64 precision and rounding to two decimal
places happens only when the result is shown. This keeps display rounding
from changing the math.

If two or more eligible bidders have the final score AGPS uses a set
six-step tie-breaking chain:

\begin{enumerate}
\item Higher final score.
\item Price.
\item Higher quality.
\item Lower delivery time.
\item Earlier submission time.
\item Lower bid identifier.
\end{enumerate}

The last criterion, bid identifier, guarantees that the ordering ends in
an order.

The evaluation functions are pure. Have no access to storage, system
time, randomness or outside input. As a result with the frozen
configuration and the same bids the engine will give the same scores and
ordering.

\begin{figure}[htbp]
\centerline{\includegraphics[width=\columnwidth]{fig2_determinism.pdf}}
\caption{Determinism contract showing configuration snapshot, canonical
JSON, SHA-256 \texttt{configHash}, pure evaluation engine, fixed
accumulation order and six-level tie-breaking cascade.}
\label{fig:determinism}
\end{figure}

\subsection{Data Provenance and Financial Sealing}
\label{subsec:provenance}

AGPS separates processing from checking the information that vendors
submit. In particular technical and quality data stay vendor-declared and
unverified. These values have a provenance record that shows their source
is \texttt{SELF\_REPORTED} and their verification status is
\texttt{UNVERIFIED}. The system has no way to change those values to
\texttt{VERIFIED}.

Financial data is also protected by a two-envelope sealing method.
Financial values cannot be seen through the API by administrators until
the tender gets to the financial-open stage. This split stops data from
being shown before the right part of the process.

The system also keeps an audit log. Every entry is added to a per-tender
SHA-256 hash chain that uses the hash and the canonical form of the
current entry. This system shows if something was changed. It does not
stop an authorized or very privileged person from changing the stored
data.

%---------------------------------------------------------------------
\section{Evaluation Methodology}
\label{sec:method}

\subsection{Eligibility Screening}

Let $R$ denote the set of enabled eligibility rules and $b$ a bid. The
screening function returns every violated rule rather than
short-circuiting on the first failure:

\begin{equation}
F(b) = \{\, r \in R \;:\; \neg\,\mathrm{satisfies}(b, r) \,\}
\label{eq:eligibility}
\end{equation}

A bid is eligible if and only if $F(b) = \emptyset$.

\subsection{Quality Derivation}

For a technical criterion $j$ with allocated points $p_j$, the earned
points depend on the declared value $v_j$ and the criterion type. For a
numeric criterion bounded by $[\,\ell_j, u_j\,]$:

\begin{equation}
t_j = \mathrm{clamp}\!\left(\frac{v_j - \ell_j}{u_j - \ell_j},\, 0,\, 1\right)
\label{eq:numeric-t}
\end{equation}

\begin{equation}
e_j =
\begin{cases}
p_j \, t_j, & \text{higher-is-better} \\[2pt]
p_j (1 - t_j), & \text{lower-is-better}
\end{cases}
\label{eq:earned}
\end{equation}

The derived quality score is $Q = \sum_j e_j$, where
$\sum_j p_j = 100$ by construction, so $Q \in [0, 100]$.

\subsection{Normalization}

Let $C$ be the cohort of eligible bids and $x_i$ the raw value of bid $i$
for a given criterion, with $x_{\min} = \min_{k \in C} x_k$ and
$x_{\max} = \max_{k \in C} x_k$. The normalized score is:

\begin{equation}
n_i =
\begin{cases}
100, & x_{\max} = x_{\min} \\[3pt]
\dfrac{x_{\min}}{x_i} \times 100, & \text{lower-is-better} \\[8pt]
\dfrac{x_i}{x_{\max}} \times 100, & \text{higher-is-better}
\end{cases}
\label{eq:normalize}
\end{equation}

The first case is a degeneracy guard: when a criterion does not
discriminate between bids it contributes uniformly rather than producing
an undefined quotient.

\subsection{Weighted Aggregation}

With integer weights $w_c$ satisfying $\sum_c w_c = 100$:

\begin{equation}
S_i = \sum_{c \in \mathcal{C}} n_{i,c} \cdot \frac{w_c}{100}
\label{eq:final}
\end{equation}

Summation follows the frozen order of $\mathcal{C}$, fixing
floating-point accumulation and yielding bit-identical results across
executions.

\subsection{Tie-Breaking}

Ranking applies a lexicographic cascade terminating in the bid
identifier, guaranteeing a total order \dots

%---------------------------------------------------------------------
\section{Implementation}
\label{sec:impl}

%---------------------------------------------------------------------
\section{Results}
\label{sec:results}

\begin{table}[htbp]
\caption{Score Breakdown for the Verification Vector}
\begin{center}
\begin{tabular}{lrrrr}
\toprule
\textbf{Criterion} & \textbf{Raw} & \textbf{Norm.} &
\textbf{Weight} & \textbf{Weighted} \\
\midrule
Price       & 800{,}000 & 92.50 & 40 & 37.00 \\
Quality     & 85        & 85.00 & 30 & 25.50 \\
Delivery    & 20        & 75.00 & 20 & 15.00 \\
Experience  & 8         & 80.00 & 10 &  8.00 \\
\midrule
\textbf{Total} & & & \textbf{100} & \textbf{85.50} \\
\bottomrule
\end{tabular}
\label{tab:golden}
\end{center}
\end{table}

\begin{table}[htbp]
\caption{Effect of Weight Configuration on Outcome (Identical Bid Data)}
\begin{center}
\begin{tabular}{lcc}
\toprule
 & \textbf{Config. A} & \textbf{Config. B} \\
\midrule
Price / Quality / Delivery / Exp. & 40/30/20/10 & 20/50/20/10 \\
Winning bidder & [NAME] & [NAME] \\
Winning score  & [VALUE] & [VALUE] \\
\bottomrule
\end{tabular}
\label{tab:flip}
\end{center}
\end{table}

%---------------------------------------------------------------------
\section{Discussion and Threats to Validity}
\label{sec:threats}

% (a) self-reported, unverified inputs
% (b) criterion verifiability taxonomy
% (c) compensatory aggregation
% (d) tamper-evident, not tamper-proof
% (e) seeded rather than live data

%---------------------------------------------------------------------
\section{Conclusion and Future Work}
\label{sec:conclusion}

%---------------------------------------------------------------------
% REFERENCES
% Option A (manual): use the thebibliography block below.
% Option B (BibTeX): delete the block, create refs.bib, and use:
%   \bibliographystyle{IEEEtran}
%   \bibliography{refs}

\begin{thebibliography}{00}

\bibitem{ref1}
% Journal article:
% A. B. Author, C. D. Author, ``Title of paper,'' \emph{Journal Name},
% vol. 1, no. 2, pp. 100--110, 2020.

\bibitem{ref2}
% Conference paper:
% A. B. Author, ``Title of paper,'' in \emph{Proc. Conf. Name}, City,
% Country, 2021, pp. 1--6.

\bibitem{ref3}
% Government document:
% Ministry of Finance, Government of India, \emph{General Financial
% Rules 2017}, New Delhi, India, 2017.

\bibitem{ref4}
% Book:
% A. B. Author, \emph{Title of Book}, 2nd ed. City: Publisher, 2019.

\bibitem{ref5}
% Online:
% Government e-Marketplace. (2024) \emph{GeM Documentation}. [Online].
% Available: https://gem.gov.in

\end{thebibliography}

\end{document}
